\section*{综合 LaTeX 正文示例}

这是一段不含 \verb|\documentclass|、导言区和宏包加载命令的正文代码，可直接放入已有文档的 \verb|\begin{document}| 与 \verb|\end{document}| 之间。内容覆盖常见数学、物理与科学写作语法，并混合使用行内公式、行间公式、分式、根式、上下标、求和、积分、极限、矩阵、向量、偏导、集合、分段思想和多种文本强调方式。

\begin{center}
\textbf{\Large 常见公式与排版语法综合练习}
\end{center}

\begin{quote}
\emph{排版的目标不是堆叠命令，而是让结构、符号和论证关系更加清楚。}
\end{quote}

\begin{itemize}
  \item \textbf{行内公式：}适合嵌入普通句子，不应过高或过长。
  \item \textbf{行间公式：}适合核心结论、长分式、矩阵和多行推导。
  \item \textbf{高亮语法：}包括粗体、斜体、下划线、边框和颜色背景。
\end{itemize}

\begin{enumerate}
  \item 先定义符号与假设；
  \item 再给出推导过程；
  \item 最后解释结果的物理或数学意义。
\end{enumerate}

\begin{description}
  \item[定义] 明确对象是什么。
  \item[定理] 陈述在什么条件下成立。
  \item[证明] 给出从条件到结论的逻辑过程。
\end{description}


\section{1. 代数与恒等变换}
代数运算是后续推导的基础。设实数变量满足 $a,b\in\mathbb{R}$，常见平方恒等式为 $(a+b)^2=a^2+2ab+b^2$，而差平方结构可写成 $a^2-b^2=(a-b)(a+b)$。在整理表达式时，应先判断是否存在公因式、完全平方或对称结构，再决定展开还是因式分解。对于含参数的式子，还要说明参数的取值范围，避免除以零或开方无意义。
\[
\sum_{k=1}^{n}k=\frac{n(n+1)}{2}
\]
\[
\prod_{k=1}^{n}k=n!
\]


\section{2. 方程与不等式}
求解方程时，等价变形必须保持解集不变。对于一元二次方程 $ax^2+bx+c=0$，判别式记为 $\Delta=b^2-4ac$，当 $\Delta\ge 0$ 时实根可写成 $x=\frac{-b\pm\sqrt{\Delta}}{2a}$。处理不等式时，需要特别注意乘除负数会改变不等号方向；含绝对值的问题通常可转化为分段讨论，并在最后取各分支解集的并集。
\[
|x-a|<r\Longleftrightarrow a-r<x<a+r
\]
\[
\frac{x-1}{x+2}\ge0\Longrightarrow x\in(-\infty,-2)\cup[1,\infty)
\]


\section{3. 函数与图像}
函数描述变量之间的对应关系。若 $y=f(x)$，定义域决定表达式在哪些位置有意义，值域则反映输出可能达到的范围。线性函数 $f(x)=kx+b$ 的斜率决定增减性，二次函数 $f(x)=ax^2+bx+c$ 的开口方向由系数 $a$ 决定。研究图像时，可以依次考察零点、对称性、单调区间、极值、渐近线以及参数变化造成的平移和伸缩。
\[
f(x+h)=f(x)+f'(x)h+o(h)
\]
\[
f^{-1}(f(x))=x
\]


\section{4. 数列与级数}
数列可看作定义在正整数上的函数。等差数列满足 $a_n=a_1+(n-1)d$，等比数列满足 $a_n=a_1q^{n-1}$，其前若干项之和常用于估计累计效应。判断无穷级数时，单看通项趋于零还不够，因为 $\lim_{n\to\infty}a_n=0$ 只是收敛的必要条件。实际分析中还会使用比较判别、比值判别、根值判别以及交错级数判别。
\[
S_n=\frac{n(a_1+a_n)}{2}
\]
\[
\sum_{n=0}^{\infty}q^n=\frac{1}{1-q},\quad |q|<1
\]


\section{5. 极限与连续}
极限刻画变量无限接近某一点时函数的趋势。经典极限 $\lim_{x\to0}\frac{\sin x}{x}=1$ 常用于三角函数的局部展开，指数函数则满足 $\lim_{n\to\infty}(1+\frac1n)^n=e$。连续性要求函数值与极限一致；若左右极限不同，则出现跳跃间断。证明极限时，可以使用夹逼、变量替换、等价无穷小或严格的定义方法。
\[
\lim_{x\to a}f(x)=L
\]
\[
\lim_{x\to0}\frac{e^x-1}{x}=1
\]


\section{6. 导数与微分}
导数反映函数的瞬时变化率。函数在一点可导意味着局部可以用线性函数近似，记作 $dy=f'(x)\,dx$。乘积求导公式为 $(uv)'=u'v+uv'$，复合函数则满足 $(f\circ g)'(x)=f'(g(x))g'(x)$。求极值时，驻点只是候选点，还需结合二阶导数、单调性或端点信息进行判断，不能把导数为零直接等同于取得极值。
\[
f'(x)=\lim_{h\to0}\frac{f(x+h)-f(x)}{h}
\]
\[
d(\ln x)=\frac{1}{x}\,dx
\]


\section{7. 积分与面积}
定积分可以理解为极限意义下的小量累加。若原函数满足 $F'(x)=f(x)$，则牛顿—莱布尼茨公式把积分计算转化为端点代入。换元积分的核心是同步替换变量与微分，分部积分则来源于乘积求导。几何应用中，面积应保持非负，因此当函数跨越横轴时需要分段取绝对值，而不能简单使用一个带符号积分。
\[
\int_a^b f(x)\,dx=F(b)-F(a)
\]
\[
\int u\,dv=uv-\int v\,du
\]


\section{8. 多元微积分}
多元函数同时依赖多个变量。对于 $z=f(x,y)$，偏导数 $\frac{\partial f}{\partial x}$ 描述固定其他变量时沿某一坐标方向的变化。梯度 $\nabla f$ 指向函数增长最快的方向，而方向导数等于梯度与单位方向向量的内积。若变量之间存在约束，可使用拉格朗日乘子法寻找候选极值，并结合边界与实际条件筛选有效结果。
\[
df=\frac{\partial f}{\partial x}dx+\frac{\partial f}{\partial y}dy
\]
\[
\nabla f=\lambda\nabla g
\]


\section{9. 线性代数}
向量与矩阵用于描述线性关系。若矩阵 $A$ 可逆，则方程组 $A\mathbf{x}=\mathbf{b}$ 的唯一解为 $\mathbf{x}=A^{-1}\mathbf{b}$。特征值问题写成 $A\mathbf{v}=\lambda\mathbf{v}$，它揭示线性变换保持方向不变的特殊向量。实际计算中通常不直接求逆矩阵，而是使用高斯消元、矩阵分解或迭代法，以获得更好的数值稳定性。
\[
\det(AB)=\det(A)\det(B)
\]
\[
A=PDP^{-1}
\]


\section{10. 概率与统计}
概率用于量化不确定性。事件的条件概率为 $P(A\mid B)=\frac{P(A\cap B)}{P(B)}$，独立性则意味着联合概率可以分解为 $P(A\cap B)=P(A)P(B)$。随机变量的期望描述平均水平，方差描述波动强度。样本统计量只能估计总体参数，因此报告结果时应同时给出样本量、误差范围和置信水平，避免只展示一个孤立的平均值。
\[
E(X)=\sum_i x_iP(X=x_i)
\]
\[
\operatorname{Var}(X)=E(X^2)-[E(X)]^2
\]


\section{11. 常微分方程}
微分方程用导数关系描述动态系统。最简单的一阶线性方程 $y'+p(x)y=q(x)$ 可通过积分因子求解；变量可分离方程则将含 $y$ 的部分与含 $x$ 的部分放到等式两侧。二阶常系数方程常借助特征方程处理。求得通解后，还必须代入初始条件或边界条件确定常数，才能得到具体问题的唯一解。
\[
y'=ky\Longrightarrow y=Ce^{kx}
\]
\[
y''+\omega^2y=0
\]


\section{12. 复数与欧拉公式}
复数把实数轴扩展为复平面。任意非零复数可写成 $z=re^{i\theta}$，其中模长为 $r=|z|$，辐角为 $\theta=\arg z$。欧拉公式 $e^{i\theta}=\cos\theta+i\sin\theta$ 将指数函数与三角函数联系起来。乘法会使模长相乘、辐角相加，因此复数特别适合处理振动、波动、电路相量以及二维旋转问题。
\[
z^n=r^ne^{in\theta}
\]
\[
\sqrt[n]{z}=r^{1/n}e^{i(\theta+2k\pi)/n}
\]


\section{13. 经典力学}
经典力学研究物体运动与受力。质点动量为 $\mathbf{p}=m\mathbf{v}$，牛顿第二定律写成 $\mathbf{F}=\frac{d\mathbf{p}}{dt}$。在保守力场中，机械能 $E=T+V$ 保持不变。处理多体系统时，质心运动与内部相对运动往往可以分离；若约束复杂，则拉格朗日方法比逐个列出约束力更加简洁，并能自然推广到广义坐标。
\[
L=T-V
\]
\[
\frac{d}{dt}\frac{\partial L}{\partial \dot q_i}-\frac{\partial L}{\partial q_i}=0
\]


\section{14. 振动与波动}
简谐振动是许多线性系统的基本近似。位移可写成 $x(t)=A\cos(\omega t+\varphi)$，速度为 $v(t)=-A\omega\sin(\omega t+\varphi)$，周期满足 $T=\frac{2\pi}{\omega}$。波动把局部振动传播到空间，波速、频率和波长之间满足确定关系。叠加原理导致干涉与驻波，而阻尼和外力会改变振幅及相位响应。
\[
v=f\lambda
\]
\[
\frac{\partial^2u}{\partial t^2}=c^2\frac{\partial^2u}{\partial x^2}
\]


\section{15. 电磁学}
电磁学用场描述电荷与电流之间的相互作用。库仑力大小与电荷乘积成正比，电场定义为单位正试探电荷所受的力，即 $\mathbf{E}=\frac{\mathbf{F}}{q}$。运动电荷还会受到磁场作用，洛伦兹力为 $\mathbf{F}=q(\mathbf{E}+\mathbf{v}\times\mathbf{B})$。麦克斯韦方程组统一了静电、稳恒磁场和电磁波，是经典场论的核心。
\[
\nabla\cdot\mathbf{E}=\frac{\rho}{\varepsilon_0}
\]
\[
\nabla\times\mathbf{B}=\mu_0\mathbf{J}+\mu_0\varepsilon_0\frac{\partial\mathbf{E}}{\partial t}
\]


\section{16. 热学与统计物理}
热力学关注宏观状态量及其约束。理想气体满足 $pV=nRT$，热力学第一定律可写成 $\Delta U=Q-W$，其中功的符号约定必须在全文保持一致。熵用于描述不可逆过程的方向，孤立系统的总熵不会自发减小。统计物理进一步把温度、压强等宏观量与微观粒子的概率分布联系起来，从而解释热力学规律的统计来源。
\[
dS=\frac{\delta Q_{\mathrm{rev}}}{T}
\]
\[
P_i=\frac{e^{-E_i/(k_BT)}}{Z}
\]


\section{17. 量子力学}
量子态由波函数或态矢量描述。波函数 $\psi(\mathbf{r},t)$ 的模平方给出概率密度，归一化要求 $\int|\psi|^2d\tau=1$。可观测量对应厄米算符，测量结果是本征值。位置与动量算符不对易，因此满足不确定关系。时间演化由薛定谔方程控制，线性叠加则允许体系同时处于多个基态或本征态的组合。
\[
i\hbar\frac{\partial}{\partial t}\psi=\hat H\psi
\]
\[
[\hat x,\hat p]=i\hbar
\]


\section{18. 相对论基础}
狭义相对论建立在光速不变与相对性原理之上。运动时钟出现时间膨胀，写作 $\Delta t=\gamma\Delta\tau$；沿运动方向的长度发生收缩，写作 $L=\frac{L_0}{\gamma}$；其中 $\gamma=\frac{1}{\sqrt{1-v^2/c^2}}$。能量与动量被统一到四维时空结构中，经典极限则在速度远小于光速时恢复。
\[
E^2=p^2c^2+m^2c^4
\]
\[
x'^2-c^2t'^2=x^2-c^2t^2
\]


\section{19. 数值计算与误差}
数值方法用有限步骤逼近精确结果。浮点运算存在舍入误差，离散化过程存在截断误差，二者可能在迭代中累积。中心差分近似 $f'(x)\approx\frac{f(x+h)-f(x-h)}{2h}$ 通常比单边差分精度更高。选择步长时不能盲目取极小值，因为过小会放大舍入误差；同时还应进行网格收敛性和稳定性检验。
\[
f'(x)=\frac{f(x+h)-f(x-h)}{2h}+O(h^2)
\]
\[
x_{n+1}=x_n-\frac{f(x_n)}{f'(x_n)}
\]


\section{20. 科学写作与排版}
科学文本不仅要给出结果，还要说明假设、符号、方法和适用范围。文中第一次出现的重要概念可用 \textbf{粗体} 强调，补充语气可用 \emph{斜体}，关键提醒可用 \underline{下划线}，短结论还可以放入 \fbox{边框}。若已加载颜色宏包，可使用 \textcolor{red}{红色警示}、\colorbox{yellow}{黄色高亮} 和 \colorbox{cyan!20}{浅色背景}。公式编号、变量字体和单位书写应保持统一，例如速度 $v=\frac{s}{t}$、功率 $P=\frac{W}{t}$、相对误差 $\varepsilon=\frac{|x-x_0|}{|x_0|}$。
\[
\mathrm{SNR}=10\log_{10}\frac{P_{\mathrm{signal}}}{P_{\mathrm{noise}}}
\]
\[
\chi^2=\sum_{i=1}^{n}\frac{(y_i-\hat y_i)^2}{\sigma_i^2}
\]



\section{21. 综合案例与表达规范}

在一篇完整的科学报告中，公式不应孤立出现，而应当与研究问题、实验条件、数据来源和结论解释相互衔接。开始写作前，可以先用一小段文字说明研究对象、观测范围和核心目标。例如，在讨论某个随时间变化的信号时，应交代采样持续时间、采样间隔、仪器量程、环境条件以及预处理方法。若数据经过平滑、插值、归一化或滤波，还应说明处理顺序和参数选择理由。这样做能够帮助读者判断结果是否可靠，也能让后续复现实验的人明确每一步操作的实际含义。

符号定义应尽量集中、清楚并保持一致。假设原始观测量记为 $x_i$，处理后的结果记为 $\hat{x}_i$，样本总数记为 $N$，则全文不宜在没有解释的情况下突然改用其他字母。下标通常表示样本编号、空间方向或不同实验组，上标可以表示幂次、转置、估计值或迭代次数。为了避免歧义，第一次出现特殊符号时应立即解释。例如，星号究竟表示共轭、最优值还是显著性，需要根据上下文明确说明。若同一个字母在不同章节承担不同含义，也应在章节开头重新定义，而不是依赖读者自行猜测。

实验步骤可以使用有序列表呈现，使操作顺序一目了然。首先检查仪器是否完成预热和校准，并记录零点漂移。其次设置采样参数，确认时间戳、通道名称和单位。随后进行空载测量，获得背景噪声水平。正式测量时应保持其他条件尽可能稳定，并为每组实验设置重复次数。完成采集后，不要立即删除异常点，而应先判断异常来自真实物理过程、操作失误还是设备故障。只有在给出明确判据后，才能决定保留、修正或剔除，并在报告中记录处理数量和原因。

数据表格应具有完整的表头、单位和有效数字。不同列的数值精度应与仪器分辨率相匹配，不能为了显示整齐而随意增加很多无意义的小数位。绘图时，横纵坐标名称、单位、图例和误差棒都应清晰可辨。若多条曲线重叠，应通过线型、标记或分图进行区分，而不是只依赖颜色。对于关键区域，可以使用局部放大图，但必须注明放大范围以及它与主图之间的对应关系。图题应说明图中展示了什么，而正文则进一步解释为什么会出现这种趋势。

误差分析至少应区分系统误差和随机误差。系统误差可能来自仪器校准、模型假设、环境偏移或方法本身的局限，通常不会通过简单增加重复次数自动消失。随机误差则表现为多次测量之间的波动，可以通过统计方法估计。样本平均值可记为 $\bar{x}$，标准差可记为字母 s。报告结果时，应说明误差量代表标准差、标准误差还是置信区间。若进行了拟合，还应展示残差分布，并检查残差是否具有明显趋势。仅给出拟合曲线看起来“很接近”并不足以证明模型正确。

讨论部分应围绕证据展开，而不是重复结果。可以先概括最重要的观察，再解释它与理论预期是否一致。如果不一致，应依次检查数据质量、参数范围、近似条件和模型结构。对于无法完全解释的现象，应明确写成限制或待解决问题，而不是用含糊语言掩盖。比较不同方法时，要保证比较基准一致，例如使用相同数据集、相同评价指标和相近计算资源。若某种方法更复杂但提升很小，也应讨论复杂度增加是否值得。

高亮语法应服务于阅读，而不是制造视觉噪声。\textbf{粗体}适合强调章节内最重要的概念，\emph{斜体}适合表示术语或轻度强调，\underline{下划线}可以用于提醒，但连续使用会降低可读性。\fbox{边框文字}适合放置很短的结论。若已有颜色支持，可以把警告写成 \textcolor{red}{需要检查的内容}，把阶段性结论写成 \colorbox{yellow}{重点结论}。颜色数量应有限，并保持相同颜色承担相同语义，避免一页中出现过多互相竞争的强调方式。

列举环境也应根据内容关系选择。没有先后顺序的内容适合使用项目符号，有明确流程的内容适合使用编号，术语及解释适合使用描述列表。若一个条目中包含较长推导，应拆分成自然段或子列表。列表前最好有一句引导语，列表后也应有一句总结，说明这些条目共同支持什么观点。层级不宜过深，否则读者会把注意力花在寻找结构上，而不是理解内容。一般情况下，两层列表已经足够表达大多数技术说明。

写结论时，应回答研究最初提出的问题。结论可以包括主要发现、适用范围、误差水平和后续改进方向，但不应加入正文从未讨论过的新证据。对于模型结果，可用决定系数 $R^2$ 或其他指标辅助评价，但任何单一指标都不能代替完整分析。对于相对误差，可说明其大小为符号 delta，同时给出绝对误差和参考值，以便读者判断误差在实际尺度上是否重要。最后还应检查引用、图表编号、公式编号和交叉引用是否正确，确保文档在修改章节顺序后仍能自动更新。

整篇文档完成后，可以从三个角度复查。第一是内容完整性：问题、方法、结果、讨论和结论是否首尾对应。第二是数学正确性：变量范围、量纲、符号、近似条件和推导步骤是否合理。第三是排版一致性：标题层级、段落间距、公式位置、标点和字体是否统一。最好再进行一次脱离作者视角的阅读，假设读者并不了解实验背景，观察哪些地方仍需要补充定义或过渡。经过这样的检查，文档不仅能够编译通过，也更容易被准确理解、验证和复现。

\section*{排版检查清单}

\begin{enumerate}
  \item 检查所有花括号、环境和数学定界符是否成对出现。
  \item 检查公式中的变量、常量、单位和算符字体是否统一。
  \item 检查长公式是否需要拆分，避免超出页面宽度。
  \item 检查列表层级是否清晰，避免连续使用过多强调颜色。
  \item 检查结论是否与前文假设一致，并说明近似成立的范围。
\end{enumerate}

\begin{center}
\fbox{\textbf{正文结束：共包含 100 个数学公式示例}}
\end{center}
